\section{研究目标与评价指标}

\subsection{数据挖掘问题定义}

本项目围绕西雅图市建筑能耗对标数据，解决以下数据挖掘问题：

\begin{table}[H]
\centering
\caption{数据挖掘问题类型与任务映射}
\label{tab:dm_problems}
\begin{tabular}{p{2.5cm}p{3cm}p{6.5cm}}
\hline
\textbf{问题类型} & \textbf{对应阶段} & \textbf{说明} \\
\hline
数据清洗与预处理 & 阶段一 & 跨年度字段审计、缺失值填补、异常值检测与处理、目标泄漏字段识别与剔除 \\
\hline
描述性统计与可视化 & 阶段二 & SiteEUI分布分析、建筑类型能耗对比、特征相关性分析、缺失模式探索 \\
\hline
监督学习—回归 & 阶段三 & 以建筑静态属性为特征，估计建筑在特定年度的合理SiteEUI基准水平；比较多种回归模型，通过滚动时间窗口交叉验证评估泛化性能 \\
\hline
模型可解释性分析 & 阶段四 & 利用SHAP方法识别对模型预测有重要贡献的特征，分析特征与预测值之间的关联模式 \\
\hline
异常检测 & 阶段五 & 基于标准化预测残差和同类统计基准，三层递进筛选异常高耗能建筑；通过持续性验证减少偶然性误判 \\
\hline
多指标综合排序 & 阶段六 & 构建多维加权评分体系，对建筑进行节能改造优先级排序，并通过多种方案对比和重采样验证排序稳定性 \\
\hline
聚类分析（辅助） & 阶段七 & 对高优先级候选建筑进行聚类画像，辅助理解候选建筑的类型结构和特征模式 \\
\hline
\end{tabular}
\end{table}

\subsection{总体目标}

以西雅图市2015--2024年建筑能耗对标数据为基础，利用建筑类型、用途、面积、层数、建造年份、区域和年度背景等预测时可获得的信息，估计建筑在特定年度的合理SiteEUI基准水平；在此基础上识别异常高耗能建筑，汇总建筑近三年持续表现，构建``数据审计 $\rightarrow$ 能耗基准预测 $\rightarrow$ 异常高耗能识别 $\rightarrow$ 建筑级汇总 $\rightarrow$ 改造优先级排序 $\rightarrow$ 辅助聚类画像''的完整分析链，为城市既有建筑节能改造初步筛选提供数据驱动的决策参考。

\subsection{各任务输入、输出与评价指标}

\begin{table}[H]
\centering
\caption{研究任务说明}
\label{tab:task_io}
\begin{tabular}{p{0.8cm}p{2.5cm}p{2.5cm}p{2.5cm}p{3.5cm}}
\hline
\textbf{序号} & \textbf{任务名称} & \textbf{输入} & \textbf{输出} & \textbf{评价指标} \\
\hline
T1 & 能耗基准回归预测 & 建筑静态属性特征矩阵、SiteEUI目标向量 & 最优回归模型及预测值 & MAE、RMSE、$R^2$；目标为显著优于全体均值基线和建筑类型中位数基线（如MAE降低10\%--20\%），不设刚性$R^2$阈值 \\
\hline
T2 & 重要预测特征识别 & 最佳模型、测试集 & SHAP特征重要性排名与可视化图表 & SHAP全局重要性排序、依赖图、瀑布图；高度相关特征分组解释；统一使用``重要预测特征''表述 \\
\hline
T3 & 异常高耗能建筑识别 & 模型预测值、实际SiteEUI、同类基准数据 & 异常建筑候选列表 & 三层递进筛选（标准化残差$>1.5$ + 同类P75基准超标 + 近三年持续出现）；报告各层筛选通过数量与重叠率 \\
\hline
T4 & 建筑级记录汇总与优先级评分 & 年度建筑记录、异常识别结果 & 建筑级优先级排名表 & 6维度加权评分；数据质量折扣修正；Top-N清单在不同权重方案下的重合率与Jaccard相似度 \\
\hline
T5 & 排序稳定性检验 & 多权重方案下的排序结果 & 稳定性分析报告 & Spearman/Kendall秩相关系数、Bootstrap入选频率、各建筑类型候选占比合理性 \\
\hline
T6 & 候选建筑聚类画像 & 高优先级候选建筑特征 & 聚类画像报告与可视化 & 轮廓系数、Calinski-Harabasz指数、画像可解释性 \\
\hline
\end{tabular}
\end{table}

\subsection{评估基线}

为验证机器学习模型的附加值，设置以下基线模型：

\begin{enumerate}
    \item \textbf{基线A（全体均值）：}以训练集SiteEUI均值作为所有建筑的预测值；
    \item \textbf{基线B（建筑类型中位数）：}按建筑类型分组，以训练集各组SiteEUI中位数作为同类建筑的预测值。
\end{enumerate}

只有当复杂模型在测试集上显著优于这些简单基线时，模型才具有实用价值。

\subsection{预测任务定义}

本项目将预测任务统一定义为：

\begin{quote}
利用建筑类型、用途、面积、层数、建造年份、区域和年度背景等\textbf{预测时可获得的信息}，估计建筑在特定年度的\textbf{合理SiteEUI基准水平}。
\end{quote}

该定义明确了项目的性质：监督学习回归与建筑能耗对标（benchmarking），而非严格的时间序列预测。除非后续明确加入上一年SiteEUI、近三年均值、同比变化率等滞后特征，否则不使用``历史能耗预测下一年能耗''的表述。
